\chapter{Αποτελέσματα και αξιολόγηση}
\label{ch:evaluation}

Το κεφάλαιο αυτό παρουσιάζει τα αποτελέσματα των μετρήσεων και τα ερμηνεύει με
βάση τους μηχανισμούς που περιγράφηκαν στα τρία προηγούμενα κεφάλαια. Όπου δεν
δηλώνεται διαφορετικά, τα διαγράμματα αφορούν ακεραίους των 32 \tl{bit} και
ομοιόμορφη κατανομή εισόδου. H επίδραση του τύπου δεδομένων και της κατανομής
εξετάζεται χωριστά στις ενότητες~\ref{sec:eval-dtypes}
και~\ref{sec:eval-distributions}.

\section{Ορθότητα}
\label{sec:eval-correctness}

Προϋπόθεση για κάθε σύγκριση χρόνου είναι η ορθότητα του αποτελέσματος. Η κύρια
σάρωση περιλαμβάνει 2700 περιπτώσεις, οι οποίες καλύπτουν και τις τρεις
αρχιτεκτονικές, τους τρεις αλγορίθμους, τους πέντε τύπους δεδομένων, μεγέθη
εισόδου από $2^3$ έως $2^{24}$ και τις τέσσερις τιμές του \tl{k}. Όλες
επαληθεύτηκαν ως προς την υλοποίηση επαλήθευσης της ενότητας~\ref{sec:meth-correctness} και
όλες πέτυχαν: ποσοστό επιτυχίας 100\% για κάθε αρχιτεκτονική και αλγόριθμο,
χωρίς καμία εξαίρεση.

Το αποτέλεσμα αυτό σημαίνει ότι κάθε χρόνος που παρουσιάζεται στη συνέχεια 
αντιστοιχεί σε υλοποίηση που παρήγαγε το σωστό σύνολο των \tl{k} στοιχείων και 
ότι οι συγκρίσεις που ακολουθούν αφορούν πράγματι τον ίδιο υπολογισμό.

\section{Χρόνος εκτέλεσης}
\label{sec:eval-time}

\subsection{Εξάρτηση από το μέγεθος εισόδου}
\label{sec:eval-time-n}

\begin{figure}[H]
    \centering
    \begin{subfigure}{0.49\textwidth}
        \centering
        \includegraphics[width=\textwidth]{int/time/time_vs_n_algo_k8}
        \caption{Αλγοριθμικό πεδίο.}
        \label{fig:eval-time-n-algo}
    \end{subfigure}
    \hfill
    \begin{subfigure}{0.49\textwidth}
        \centering
        \includegraphics[width=\textwidth]{int/time/time_vs_n_e2e_k8}
        \caption{Πεδίο \tl{end-to-end}.}
        \label{fig:eval-time-n-e2e}
    \end{subfigure}
    \caption[Χρόνος συναρτήσει του μεγέθους, στα δύο πεδία μέτρησης]{Χρόνος συναρτήσει του μεγέθους εισόδου, για \tl{k} ίσο με 8, στα
    δύο πεδία μέτρησης. Η κατάταξη των αρχιτεκτονικών ανατρέπεται πλήρως
    μεταξύ των δύο.}
    \label{fig:eval-time-n-k8}
\end{figure}

Στο αριστερό σχήμα, που αφορά το αλγοριθμικό πεδίο, η μονάδα γραφικών
είναι ταχύτερη στη μικρότερη είσοδο,
ολοκληρώνοντας σε \qty{0.010}{\milli\second} έναντι \qty{0.022}{\milli\second}
του επεξεργαστή, και πάλι στη μεγαλύτερη, σε \qty{0.43}{\milli\second} έναντι
\qty{1.23}{\milli\second}· σχεδόν σε ολόκληρη την ενδιάμεση περιοχή όμως ο
επεξεργαστής είναι ο ταχύτερος από τους δύο. Το ακριβές μέγεθος στο οποίο οι δύο εναλλάσσουν θέσεις είναι
από μόνο του ασταθές μεταξύ επαναλαμβανόμενων εκτελέσεων της ίδιας σάρωσης:
σε αυτούς τους μικρούς απόλυτους χρόνους, αμφότερες οι αρχιτεκτονικές
βρίσκονται αρκετά κοντά στην ανάλυση του χρονομέτρου ώστε η κατάταξη εκεί να
αντικατοπτρίζει θόρυβο μέτρησης και όχι σταθερή αρχιτεκτονική διαφορά.

Η μονάδα νευρωνικής επεξεργασίας εμφανίζει σαφές κατώφλι στα μικρά μεγέθη. Οι
καμπύλες της παραμένουν εντελώς επίπεδες γύρω
στο μισό χιλιοστό του δευτερολέπτου μέχρι περίπου $2^{16}$ στοιχεία: σε όλη
αυτή την περιοχή ο χρόνος δεν εξαρτάται καθόλου από το μέγεθος της εισόδου, διότι
κυριαρχείται από το σταθερό κόστος αποστολής και αναμονής. Ο επεξεργαστής
παρουσιάζει ανάλογο, αλλά πολύ χαμηλότερο κατώφλι, της τάξης των πέντε
εκατοστών του χιλιοστού.

Ένα σημείο ωστόσο ξεχωρίζει από αυτή την εικόνα: στο μικρότερο μέγεθος
εισόδου, όπου \tl{n} και \tl{k} ταυτίζονται, το \tl{map-reduce} της μονάδας
νευρωνικής επεξεργασίας εμφανίζει χρόνο αισθητά κάτω από το κατώφλι που
μόλις περιγράφηκε. Ο λόγος βρίσκεται στη δειγματοληψία της
ενότητας~\ref{sec:npu-mapreduce}: όταν \tl{n}=\tl{k}=8, το δείγμα που
τροφοδοτεί τον σωρό καλύπτει ολόκληρη την είσοδο, οπότε καμία εργασία δεν
αποστέλλεται ποτέ στα πλακίδια. Το σημείο αυτό μετρά επομένως αποκλειστικά
την πλευρά του \tl{host} και δεν αντιπροσωπεύει πραγματική συμμετοχή της
διάταξης.

Η εικόνα αυτή όμως ανατρέπεται πλήρως στο δεξιό σχήμα, όταν η μέτρηση
περιλαμβάνει και τη μεταφορά των δεδομένων. Στο πεδίο \tl{end-to-end} ο
επεξεργαστής είναι ταχύτερος
από κάθε άλλη αρχιτεκτονική σε \emph{όλα} τα μεγέθη που εξετάστηκαν,
συμπεριλαμβανομένου του μεγαλύτερου: τα \qty{0.43}{\milli\second} του πυρήνα της
μονάδας γραφικών γίνονται \qty{5.69}{\milli\second} συνολικά, ενώ ο επεξεργαστής
παραμένει στα \qty{1.23}{\milli\second}, καθώς για αυτόν τα δύο πεδία
ταυτίζονται.

Το αλγοριθμικό πεδίο μετρά πόσο καλά εκτελεί η κάθε συσκευή τον
υπολογισμό· το συνολικό μετρά αν αξίζει να της ανατεθεί. Οι δύο απαντήσεις δεν
συμπίπτουν, και η ενότητα~\ref{sec:eval-energy} δείχνει ότι η δεύτερη είναι
εκείνη που καθορίζει και την ενεργειακή εικόνα. Η μονάδα νευρωνικής
επεξεργασίας απαιτεί \qty{9.97}{\milli\second} στη μεγαλύτερη είσοδο,
περισσότερο από κάθε άλλη αρχιτεκτονική και στα δύο πεδία, για τους λόγους
της ενότητας~\ref{sec:npu-bottlenecks}.

\begin{figure}[H]
    \centering
    \begin{subfigure}{0.49\textwidth}
        \centering
        \includegraphics[width=\textwidth]{int/time/time_vs_n_backend_compare_bitonic_k8}
        \caption{Διτονική ταξινόμηση.}
        \label{fig:eval-backend-bitonic}
    \end{subfigure}
    \hfill
    \begin{subfigure}{0.49\textwidth}
        \centering
        \includegraphics[width=\textwidth]{int/time/time_vs_n_backend_compare_map_reduce_k8}
        \caption{\tl{Map-reduce}.}
        \label{fig:eval-backend-mapreduce}
    \end{subfigure}
    \caption[Κάθε αλγόριθμος στις τρεις αρχιτεκτονικές]{Κάθε αλγόριθμος στις τρεις αρχιτεκτονικές, για \tl{k} ίσο με 8.}
    \label{fig:eval-backend-compare}
\end{figure}

Η σύγκριση της ίδιας διτονικής ταξινόμησης στις τρεις αρχιτεκτονικές είναι
διαφωτιστική, διότι εκτελείται το ίδιο χρονοδιάγραμμα επιπέδων και η διαφορά
οφείλεται αποκλειστικά στην υλοποίηση. Οι δύο αρχιτεκτονικές γενικού σκοπού
εναλλάσσουν θέσεις κατά μήκος του εύρους: ο επεξεργαστής είναι ταχύτερος από
τη μονάδα γραφικών στο μεγαλύτερο μέρος των μικρών έως μεσαίων μεγεθών, και η
μονάδα γραφικών προηγείται μόνο στα μεγαλύτερα μεγέθη, φτάνοντας παράγοντα έξι στο
μεγαλύτερο μέγεθος που εξετάστηκε, συνεπές με τη διαφορά του διαθέσιμου
εύρους ζώνης. Όπως και στην παραπάνω σύγκριση, το ακριβές μέγεθος στο οποίο
οι δύο εναλλάσσουν θέσεις μετακινείται μεταξύ επαναλαμβανόμενων εκτελέσεων
της ίδιας σάρωσης και δεν πρέπει να διαβάζεται ως σταθερό όριο. Η μονάδα
νευρωνικής επεξεργασίας βρίσκεται κοντά στον επεξεργαστή σε αυτά τα
μεγαλύτερα μεγέθη, κάτι που είναι συνεπές με το εύρημα ότι ο χρόνος της
δαπανάται σχεδόν εξ ολοκλήρου στον \tl{host}.

Το \tl{map-reduce} δεν επιδέχεται καταρχήν την ίδια σύγκριση, καθώς κάθε
αρχιτεκτονική διατηρεί και συγχωνεύει τους τοπικούς σωρούς της με θεμελιωδώς
διαφορετικό τρόπο (κεφάλαια~\ref{ch:cpu} έως~\ref{ch:npu})· το δεξιό σχήμα
δείχνει αντ' αυτού πόσο διαφορετικά εκδιπλώνεται η ίδια αλγοριθμική ιδέα όταν
υλοποιηθεί τρεις φορές. Η μονάδα γραφικών είναι η ταχύτερη των τριών σχεδόν
σε ολόκληρο το εύρος, με τον επεξεργαστή να ακολουθεί σε περίπου σταθερή
απόσταση, εκτός από τα μικρά μεγέθη όπου ισχύει η ίδια επιφύλαξη περί
θορύβου μέτρησης όπως παραπάνω· στη μεγαλύτερη είσοδο η μονάδα γραφικών
παραμένει η ταχύτερη, όπως και παραπάνω, χωρίς όμως την ενδιάμεση εναλλαγή
θέσεων που εμφανίζει η διτονική ταξινόμηση. Η μονάδα νευρωνικής επεξεργασίας, σε έντονη αντίθεση
με τη συμπεριφορά της στη διτονική ταξινόμηση, δεν βρίσκεται εδώ κοντά στον
επεξεργαστή: παραμένει περίπου τάξη μεγέθους βραδύτερη και από τις δύο
αρχιτεκτονικές σε κάθε μέγεθος, φτάνοντας τον ήδη αναφερθέντα χρόνο της στη
μεγαλύτερη είσοδο. Ο λόγος προκύπτει από την ενότητα~\ref{sec:npu-bottlenecks}: ο
ρόλος της διάταξης στο \tl{map-reduce} είναι μόνο ένα φίλτρο ανά τμήμα και
όχι κοινό χρονοδιάγραμμα δικτύου, οπότε καμία από την ομοιότητα δομής που
κρατά την καμπύλη της διτονικής ταξινόμησης κοντά στον επεξεργαστή δεν
μεταφέρεται σε αυτόν τον αλγόριθμο.

\subsection{Εξάρτηση από την παράμετρο \tl{k}}
\label{sec:eval-time-k}

Η εξάρτηση από το \tl{k} είναι το σημείο όπου η κατάταξη των αλγορίθμων
ανατρέπεται, και αποτελεί ίσως το σημαντικότερο εύρημα της εργασίας.

\begin{figure}[H]
    \centering
    \begin{subfigure}{0.49\textwidth}
        \centering
        \includegraphics[width=\textwidth]{int/time/time_vs_n_k_colored_cpu_bitonic}
        \caption{Διτονική ταξινόμηση, \tl{CPU}.}
    \end{subfigure}
    \hfill
    \begin{subfigure}{0.49\textwidth}
        \centering
        \includegraphics[width=\textwidth]{int/time/time_vs_n_k_colored_gpu_map_reduce}
        \caption{\tl{Map-reduce}, \tl{GPU}.}
    \end{subfigure}
    \caption[Χρόνος συναρτήσει του μεγέθους, ανά τιμή του \tl{k}]{Χρόνος συναρτήσει του μεγέθους εισόδου, με χωριστή καμπύλη ανά τιμή
    του \tl{k}. Οι καμπύλες της διτονικής ταξινόμησης παραμένουν κοντά μεταξύ
    τους, ενώ εκείνες του \tl{map-reduce} αποκλίνουν κατά τάξεις μεγέθους.}
    \label{fig:eval-k-colored}
\end{figure}

Η αντιπαράθεση του σχήματος~\ref{fig:eval-k-colored} αποτυπώνει το φαινόμενο
άμεσα. Για τη διτονική ταξινόμηση οι καμπύλες των τεσσάρων τιμών του \tl{k}
σχεδόν ταυτίζονται, καθώς το \tl{k} επηρεάζει μόνο το μέγεθος των ταξινομημένων
μπλοκ. Για το \tl{map-reduce} στη μονάδα γραφικών οι ίδιες καμπύλες απέχουν
μεταξύ τους τρεις τάξεις μεγέθους.

\begin{table}[H]
    \centering
    \begin{tabular}{lrrrr}
        \hline
        & \multicolumn{2}{c}{\tl{k}=8} & \multicolumn{2}{c}{\tl{k}=131072} \\
        Υλοποίηση & \tl{ms} & θέση & \tl{ms} & θέση \\
        \hline
        \tl{CPU} διτονική & 8,01 & 4 & 43,33 & 2 \\
        \tl{CPU map-reduce} & 1,23 & 2 & 60,37 & 5 \\
        \tl{GPU} διτονική & 1,34 & 3 & 8,22 & 1 \\
        \tl{GPU map-reduce} & 0,43 & 1 & 1191,37 & 6 \\
        \tl{NPU} διτονική & 8,17 & 5 & 56,70 & 3 \\
        \tl{NPU map-reduce} & 9,97 & 6 & 59,69 & 4 \\
        \hline
    \end{tabular}
    \caption[Αλγοριθμικός χρόνος για τα δύο άκρα του \tl{k}]{Αλγοριθμικός χρόνος για τα δύο άκρα του \tl{k}, σε είσοδο $2^{24}$
    στοιχείων. Η κατάταξη αφορά τις έξι υλοποιήσεις της εργασίας και δεν
    περιλαμβάνει τις υλοποιήσεις αναφοράς.}
    \label{tab:eval-k-inversion}
\end{table}

Η ταχύτερη υλοποίηση για μικρό \tl{k} γίνεται η βραδύτερη για μεγάλο, με διαφορά
τριών τάξεων μεγέθους. Το \tl{map-reduce} στη μονάδα γραφικών επιβραδύνεται κατά
παράγοντα 2800 μεταξύ των δύο άκρων, ενώ η διτονική ταξινόμηση στην ίδια
αρχιτεκτονική μόλις κατά έξι. Για μικρό \tl{k} οι δύο υλοποιήσεις του
επεξεργαστή κατατάσσονται πλέον υψηλότερα σε σχέση με τη θέση τους στο μεγάλο
\tl{k}, με το \tl{map-reduce} να προσπερνά τη διτονική της μονάδας γραφικών
και τη διτονική να προσπερνά εκείνη της μονάδας νευρωνικής επεξεργασίας, όμως
η ίδια η ανατροπή παραμένει ανεπηρέαστη: η φθηνότερη υλοποίηση για \tl{k}
ίσο με 8 δεν είναι ποτέ η φθηνότερη για \tl{k} ίσο με 131072.

Η ανατροπή αυτή δεν είναι ιδιαιτερότητα της παρούσας υλοποίησης. Ανεξάρτητες
μελέτες σε μονάδες γραφικών δύο διαφορετικών κατασκευαστών καταγράφουν την ίδια
εξάρτηση από το \tl{k} και την αντιμετωπίζουν με προσαρμοστική εναλλαγή
αλγορίθμου~\cite{zhang2023gputopk,lin2026adaptivetopk}· η αύξηση του χρόνου των
μεθόδων μερικής ταξινόμησης με το \tl{k} αποδίδεται εκεί ρητά στην πολυπλοκότητα
$O(\log^2 k)$ του υποκείμενου διτονικού δικτύου.

Η αιτία δόθηκε στην ενότητα~\ref{sec:gpu-mapreduce-limit}: ο σωρός φιλτράρει μόνο
αφού γεμίσει, οπότε το πλήθος των νημάτων που συνεισφέρουν περιορίζεται στον
λόγο $n/k$. Για τα μεγέθη του πίνακα~\ref{tab:eval-k-inversion} ο λόγος αυτός
είναι 128, δηλαδή η συσκευή απασχολεί ουσιαστικά εκατόν είκοσι οκτώ νήματα από
τα χιλιάδες που διαθέτει. Ο επεξεργαστής, με μόλις οκτώ νήματα, δεν προσκρούει
ποτέ στο ίδιο όριο και για τον λόγο αυτό η επιβράδυνσή του είναι σαφώς ηπιότερη.

\begin{figure}[H]
    \centering
    \includegraphics[width=0.9\textwidth]{int/time/time_vs_n_algo_compare_gpu_k131072}
    \caption{Οι τρεις αλγόριθμοι στη μονάδα γραφικών για \tl{k} ίσο με 131072.}
    \label{fig:eval-algo-compare-gpu}
\end{figure}

Το πρακτικό συμπέρασμα είναι ότι δεν υπάρχει καθολικά προτιμότερος αλγόριθμος.
Η επιλογή μεταξύ δικτύου ταξινόμησης και επιλογής με σωρό εξαρτάται από τον λόγο
$n/k$ και όχι από την αρχιτεκτονική.

\section{Σύγκριση αρχιτεκτονικών}
\label{sec:eval-backends}

\begin{figure}[H]
    \centering
    \begin{subfigure}{0.49\textwidth}
        \centering
        \includegraphics[width=\textwidth]{int/time/speedup_vs_gt_cpu_k8}
        \caption{\tl{CPU}, \tl{k}=8.}
    \end{subfigure}
    \hfill
    \begin{subfigure}{0.49\textwidth}
        \centering
        \includegraphics[width=\textwidth]{int/time/speedup_vs_gt_cpu_k131072}
        \caption{\tl{CPU}, \tl{k}=131072.}
    \end{subfigure}

    \vspace{0.25em}

    \begin{subfigure}{0.49\textwidth}
        \centering
        \includegraphics[width=\textwidth]{int/time/speedup_vs_gt_gpu_k8}
        \caption{\tl{GPU}, \tl{k}=8.}
    \end{subfigure}
    \hfill
    \begin{subfigure}{0.49\textwidth}
        \centering
        \includegraphics[width=\textwidth]{int/time/speedup_vs_gt_gpu_k131072}
        \caption{\tl{GPU}, \tl{k}=131072.}
    \end{subfigure}

    \vspace{0.25em}

    \begin{subfigure}{0.49\textwidth}
        \centering
        \includegraphics[width=\textwidth]{int/time/speedup_vs_gt_npu_k8}
        \caption{\tl{NPU}, \tl{k}=8.}
    \end{subfigure}
    \hfill
    \begin{subfigure}{0.49\textwidth}
        \centering
        \includegraphics[width=\textwidth]{int/time/speedup_vs_gt_npu_k131072}
        \caption{\tl{NPU}, \tl{k}=131072.}
    \end{subfigure}
    \caption[Επιτάχυνση έναντι της υλοποίησης αναφοράς]{Επιτάχυνση έναντι της υλοποίησης αναφοράς της ίδιας αρχιτεκτονικής.
    Αριστερή στήλη \tl{k} ίσο με 8, δεξιά στήλη \tl{k} ίσο με 131072, μία γραμμή
    ανά αρχιτεκτονική. Τιμές πάνω από τη μονάδα σημαίνουν ότι η υλοποίηση της
    εργασίας είναι ταχύτερη από τη βιβλιοθήκη.}
    \label{fig:eval-speedup}
\end{figure}

Η σύγκριση κάθε υλοποίησης με το μέτρο σύγκρισης της ίδιας αρχιτεκτονικής δίνει
την πιο ειλικρινή εικόνα, διότι εξουδετερώνει τις διαφορές του υλικού.

Στον επεξεργαστή, το \tl{map-reduce} είναι ταχύτερο της βιβλιοθήκης κατά
παράγοντα τρία για μικρό \tl{k}. Πρέπει όμως να υπενθυμιστεί ο περιορισμός της
ενότητας~\ref{sec:cpu-reference}: η βιβλιοθήκη είναι μονονηματική ενώ η
υλοποίηση χρησιμοποιεί οκτώ νήματα, οπότε ο παράγοντας αυτός δεν είναι καθαρά
αλγοριθμικός.

Η μεγαλύτερη επιτάχυνση ολόκληρου του σχήματος εμφανίζεται ωστόσο αλλού, και
είναι συνέπεια της ανατροπής της ενότητας~\ref{sec:eval-time-k}. Για το
μεγάλο \tl{k}, όπου η επιλογή με σωρό καταρρέει, η διτονική ταξινόμηση στον
επεξεργαστή φτάνει να είναι περίπου έξι φορές ταχύτερη από τη βιβλιοθήκη. Ο
λόγος είναι ότι η βιβλιοθήκη στηρίζεται και αυτή σε σωρό μεγέθους \tl{k} και
υφίσταται την ίδια επιβάρυνση, ενώ το δίκτυο ταξινόμησης παραμένει εν πολλοίς
ανεπηρέαστο. Το πλεονέκτημα εξασθενεί καθώς η είσοδος μεγαλώνει περαιτέρω και η
κίνηση δεδομένων του δικτύου αρχίζει να κυριαρχεί, περιοριζόμενο σε λίγο πάνω
από τη μονάδα στη μεγαλύτερη είσοδο.

Στη μονάδα γραφικών η σύγκριση είναι αυστηρή, καθώς η βιβλιοθήκη εκτελείται στην
ίδια συσκευή με την ίδια παραλληλία. Εκεί το \tl{map-reduce} υπερτερεί οριακά
για \tl{k} ίσο με 8, ενώ για μεγάλο \tl{k} η βιβλιοθήκη είναι ταχύτερη κατά
περισσότερο από χίλιες φορές. Η επιλογή με βάση την αναπαράσταση δεν έχει σημείο
κατάρρευσης ως προς το \tl{k}, όπως εξηγήθηκε στην
ενότητα~\ref{sec:gpu-reference-cub}.

Στη μονάδα νευρωνικής επεξεργασίας, για τους ακέραιους τύπους στους οποίους
αναφέρεται το σχήμα, καμία υλοποίηση δεν ξεπερνά τη βιβλιοθήκη σε καμία περιοχή
παραμέτρων. Το αποτέλεσμα είναι αναμενόμενο και όχι απογοητευτικό: όπως δείχθηκε
στην ενότητα~\ref{sec:npu-bottlenecks}, ο ίδιος επεξεργαστής που θα εκτελούσε τη
βιβλιοθήκη καταναλώνεται ούτως ή άλλως στην προετοιμασία των δεδομένων για τη
μονάδα.

Ο περιορισμός «για τους ακέραιους τύπους» δεν είναι διατύπωση ασφαλείας. Για
τους τύπους κινητής υποδιαστολής, και ιδίως για τη μισή ακρίβεια, η εικόνα
ανατρέπεται και η μονάδα νευρωνικής επεξεργασίας πράγματι ξεπερνά τη
βιβλιοθήκη· η εξήγηση δίνεται στην ενότητα~\ref{sec:eval-dtypes}, αφού πρώτα
παρουσιαστεί η συμπεριφορά των τύπων.

\section{Ανάλυση με το μοντέλο \tl{roofline}}
\label{sec:eval-roofline}

\begin{figure}[H]
    \centering
    \begin{subfigure}{\textwidth}
        \centering
        \includegraphics[width=0.86\textwidth]{int/memory/roofline_kernels_k8}
        \caption{\tl{k} ίσο με 8.}
        \label{fig:eval-roofline-k8}
    \end{subfigure}

    \vspace{0.3em}

    \begin{subfigure}{\textwidth}
        \centering
        \includegraphics[width=0.86\textwidth]{int/memory/roofline_kernels_k131072}
        \caption{\tl{k} ίσο με 131072.}
        \label{fig:eval-roofline-k131072}
    \end{subfigure}
    \caption[Τοποθέτηση των εκτελέσεων στα μετρημένα ταβάνια]{Τοποθέτηση των εκτελέσεων απέναντι στα μετρημένα ταβάνια, για τη
    μεγαλύτερη είσοδο. Περιορίζεται στο μέγιστο μέγεθος, όπου το σύνολο εργασίας
    ξεπερνά την κρυφή μνήμη και το ταβάνι της κύριας μνήμης είναι το σχετικό.}
    \label{fig:eval-roofline}
\end{figure}

Το διάγραμμα τοποθετεί κάθε εκτέλεση στο επίπεδο υπολογιστικής έντασης και
επιτευχθέντος ρυθμού. Όλες οι υλοποιήσεις συγκεντρώνονται στο αριστερό τμήμα,
δηλαδή σε πολύ χαμηλή υπολογιστική ένταση, και καμία δεν πλησιάζει το οριζόντιο
τμήμα της στέγης. Καμία υλοποίηση της εργασίας, σε καμία αρχιτεκτονική, δεν
περιορίζεται από την υπολογιστική ικανότητα. Το εύρημα συμφωνεί με ανεξάρτητη
μέτρηση σε επιταχυντές κέντρου δεδομένων, όπου ο ταχύτερος γνωστός πυρήνας
επιλογής αξιοποιεί πάνω από το ενενήντα τοις εκατό του διαθέσιμου εύρους ζώνης
έναντι περίπου του μισού της υπολογιστικής ικανότητας, και όπου ο χρόνος
κλιμακώνεται με το εύρος ζώνης της κάθε συσκευής~\cite{zhang2023gputopk}.

Η αντιπαράθεση των δύο τιμών του \tl{k} προσθέτει δύο παρατηρήσεις. Και στις δύο
περιπτώσεις η διτονική ταξινόμηση στη μονάδα γραφικών βρίσκεται πρακτικά πάνω
στο κεκλιμένο τμήμα της δικής της στέγης, κάτι που επιβεβαιώνει οπτικά το εύρημα
της ενότητας~\ref{sec:gpu-bottlenecks}: η υλοποίηση αξιοποιεί το διαθέσιμο εύρος
ζώνης σχεδόν πλήρως και περιορίζεται αποκλειστικά από τον όγκο που υποχρεούται
να διακινήσει.

Η δεύτερη παρατήρηση αφορά το \tl{map-reduce} της ίδιας αρχιτεκτονικής. Για
\tl{k} ίσο με 8 βρίσκεται κοντά στα υπόλοιπα σημεία, ενώ για τη μεγάλη τιμή
καταρρέει κατά τρεις τάξεις μεγέθους και απομακρύνεται εξίσου από αμφότερα τα
ταβάνια. Η θέση αυτή, μακριά τόσο από το κεκλιμένο όσο και από το οριζόντιο
τμήμα, είναι η υπογραφή ενός περιορισμού που δεν είναι ούτε η μνήμη ούτε ο
υπολογισμός, και αντιστοιχεί στην κατάρρευση της παραλληλίας της
ενότητας~\ref{sec:gpu-mapreduce-limit}.

Το εύρημα αυτό είναι δομικό και όχι συγκυριακό. Η επιλογή \tl{top-k} εκτελεί
ουσιαστικά μία σύγκριση ανά στοιχείο και δεν επαναχρησιμοποιεί σχεδόν τίποτα,
οπότε η υπολογιστική της ένταση είναι χαμηλή εξ ορισμού.

\begin{figure}[H]
    \centering
    \includegraphics[width=0.85\textwidth]{machine/cache_ladder}
    \caption{Μετρημένο εύρος ζώνης ανάγνωσης συναρτήσει του μεγέθους του συνόλου
    εργασίας, για τις τρεις αρχιτεκτονικές.}
    \label{fig:eval-cache-ladder}
\end{figure}

Το ταβάνι με το οποίο συγκρίνεται κάθε εκτέλεση δεν είναι σταθερό αλλά
συνάρτηση του συνόλου εργασίας της, σύμφωνα με την
ενότητα~\ref{sec:meth-roofline}. Το σχήμα~\ref{fig:eval-cache-ladder} δείχνει τη
μετρημένη κλίμακα: τα σκαλοπάτια αντιστοιχούν στις βαθμίδες της ιεραρχίας
μνήμης, και η πτώση προς το δεξιό άκρο σηματοδοτεί το σημείο όπου το σύνολο
εργασίας παύει να χωρά στην τελευταία βαθμίδα.

\subsection{Αξιοποίηση εύρους ζώνης μνήμης}
\label{sec:eval-bandwidth}

\begin{figure}[H]
    \centering
    \begin{subfigure}{\textwidth}
        \centering
        \includegraphics[width=0.62\textwidth]{int/memory/cpu_memory_bandwidth_vs_n_k8}
        \caption{\tl{CPU}.}
    \end{subfigure}

    \vspace{0.15em}

    \begin{subfigure}{\textwidth}
        \centering
        \includegraphics[width=0.62\textwidth]{int/memory/gpu_memory_bandwidth_vs_n_k8}
        \caption{\tl{GPU}.}
    \end{subfigure}

    \vspace{0.15em}

    \begin{subfigure}{\textwidth}
        \centering
        \includegraphics[width=0.62\textwidth]{int/memory/npu_memory_bandwidth_vs_n_k8}
        \caption{\tl{NPU}.}
    \end{subfigure}
    \caption[Επιτευχθέν εύρος ζώνης ανά αρχιτεκτονική]{Επιτευχθέν εύρος ζώνης συναρτήσει του μεγέθους εισόδου, μία
    αρχιτεκτονική ανά σχήμα, με το μετρημένο ταβάνι της σχεδιασμένο με
    διακεκομμένη γραμμή. Η απόσταση κάθε καμπύλης από τη διακεκομμένη είναι το
    αναξιοποίητο περιθώριο.}
    \label{fig:eval-bandwidth}
\end{figure}

Η αξιοποίηση του ταβανιού μνήμης αυξάνεται με το μέγεθος της εισόδου και
σταθεροποιείται όταν το σύνολο εργασίας ξεπεράσει την κρυφή μνήμη. Στη
μεγαλύτερη είσοδο, η διτονική ταξινόμηση στον επεξεργαστή φτάνει περίπου τα
τέσσερα πέμπτα του ταβανιού, αισθητά μπροστά από τα τρία πέμπτα της μονάδας
γραφικών, ενώ η αξιοποίηση του υπολογιστικού ταβανιού παραμένει κάτω του
πέντε τοις εκατό και στις δύο.

Η αντιπαράθεση της διτονικής ταξινόμησης με τη βιβλιοθήκη στη μονάδα γραφικών
αξίζει ιδιαίτερη προσοχή, διότι δείχνει ότι η εγγύτητα στο ταβάνι δεν ισοδυναμεί
με ταχύτητα. Η διτονική ταξινόμηση αξιοποιεί σαφώς μεγαλύτερο ποσοστό του
διαθέσιμου εύρους ζώνης από τη βιβλιοθήκη, και παρ' όλα αυτά είναι βραδύτερη,
επειδή διακινεί επταπλάσιο όγκο δεδομένων. Η αξιοποίηση μετρά την ποιότητα της
υλοποίησης· ο συνολικός όγκος μετρά την ποιότητα του αλγορίθμου.

\section{Ενεργειακή αποτίμηση}
\label{sec:eval-energy}

\begin{figure}[H]
    \centering
    \includegraphics[width=0.9\textwidth]{int/time/time_vs_n_metric_compare_gpu_k8}
    \caption[Αλγοριθμικός και συνολικός χρόνος στη μονάδα γραφικών]{Αλγοριθμικός και συνολικός χρόνος στη μονάδα γραφικών, για k ίσο με 8. Η απόσταση
    των δύο καμπυλών κάθε αλγορίθμου είναι το κόστος των μεταφορών και της
    έναρξης των πυρήνων.}
    \label{fig:eval-metric-compare}
\end{figure}

Πριν από την παρουσίαση των ενεργειακών μετρήσεων είναι χρήσιμο το
σχήμα~\ref{fig:eval-metric-compare}, το οποίο δείχνει γιατί τα δύο πεδία
μέτρησης αποκλίνουν τόσο στη μονάδα γραφικών. Η καμπύλη του συνολικού χρόνου
παραμένει σταθερή κοντά στα τρία δέκατα του χιλιοστού του δευτερολέπτου μέχρι
περίπου $2^{16}$ στοιχεία: σε όλη αυτή την περιοχή ο πυρήνας είναι ουσιαστικά
αόρατος και ο χρόνος αναλώνεται εξ ολοκλήρου σε μεταφορές και εκκινήσεις. Ακόμη
και στη μεγαλύτερη είσοδο, ο συνολικός χρόνος παραμένει περίπου τριπλάσιος του
αλγοριθμικού.

Η ενεργειακή εικόνα διαφέρει ουσιωδώς από την εικόνα των χρόνων, και η διαφορά
οφείλεται στο συμμετρικό πεδίο μέτρησης της ενότητας~\ref{sec:meth-energy-scope}:
κάθε αρχιτεκτονική χρεώνεται το συνολικό κόστος επίλυσης, συμπεριλαμβανομένου
του \tl{host}.

Στη μεγαλύτερη είσοδο και για μικρό \tl{k}, το \tl{map-reduce} στον επεξεργαστή
καταναλώνει περίπου \qty{19}{\milli\joule} ανά εκτέλεση και είναι η οικονομικότερη
υλοποίηση. Οι αντίστοιχες υλοποιήσεις στη μονάδα γραφικών καταναλώνουν από
\qty{108}{\milli\joule} έως \qty{128}{\milli\joule}, δηλαδή έως επτά φορές
περισσότερο, παρότι είναι ταχύτερες στο αλγοριθμικό πεδίο. Η ακριβότερη
υλοποίηση όλων είναι η διτονική ταξινόμηση στη μονάδα γραφικών, ενώ η αντίστοιχη
στον επεξεργαστή κοστίζει \qty{118}{\milli\joule}.

Η φαινομενική αντίφαση εξηγείται από την ενότητα~\ref{sec:gpu-memory}. Το
αλγοριθμικό πεδίο της μονάδας γραφικών περιλαμβάνει μόνο τον πυρήνα, ενώ το
συνολικό περιλαμβάνει τις μεταφορές, οι οποίες καταναλώνουν χρόνο και ισχύ και
από τις δύο πλευρές του διαύλου. Μια συσκευή που είναι τέσσερις φορές ταχύτερη
στον πυρήνα, αλλά χρειάζεται δύο φορές τον χρόνο του πυρήνα μόνο για να λάβει τα
δεδομένα, και καταναλώνει ταυτόχρονα ισχύ πλακέτας, δεν είναι οικονομικότερη.

\begin{figure}[H]
    \centering
    \begin{subfigure}{0.49\textwidth}
        \centering
        \includegraphics[width=\textwidth]{int/energy/energy_by_backend_net_k8}
        \caption{Καθαρή ενέργεια.}
        \label{fig:eval-energy}
    \end{subfigure}
    \hfill
    \begin{subfigure}{0.49\textwidth}
        \centering
        \includegraphics[width=\textwidth]{int/energy/power_by_backend_net_k8}
        \caption{Καθαρή ισχύς.}
        \label{fig:eval-power}
    \end{subfigure}
    \caption[Καθαρή ενέργεια και ισχύς ανά αρχιτεκτονική και αλγόριθμο]{Καθαρή ενέργεια και ισχύς ανά αρχιτεκτονική και αλγόριθμο, για \tl{k} ίσο με 8.}
    \label{fig:eval-energy-power}
\end{figure}

Η αντιπαράθεση των σχημάτων~\ref{fig:eval-energy} και~\ref{fig:eval-power}
αναδεικνύει μια διάκριση που είναι εύκολο να συγχέεται, και η οποία αφορά
ιδιαίτερα τη μονάδα νευρωνικής επεξεργασίας. Η μονάδα αυτή έχει τη
\emph{χαμηλότερη} κατανάλωση ισχύος από όλες τις υλοποιήσεις, με το
\tl{map-reduce} να μετριέται στα \qty{7.9}{\watt} έναντι \qty{15.6}{\watt} του
αντίστοιχου στον επεξεργαστή και \qty{18.6}{\watt} της μονάδας γραφικών.

Παρ' όλα αυτά καταναλώνει πάνω από τετραπλάσια ενέργεια από το \tl{map-reduce} του
επεξεργαστή, \qty{0.080}{\joule} έναντι \qty{0.019}{\joule}. Ο λόγος είναι ότι η
ενέργεια είναι το γινόμενο της ισχύος επί τον χρόνο, και η μονάδα χρειάζεται
περίπου οκταπλάσιο χρόνο για την ίδια εργασία. Το πλεονέκτημα του σαράντα εννέα τοις
εκατό στην ισχύ εξανεμίζεται από το μειονέκτημα του οκταπλασίου στη διάρκεια.

Η παρατήρηση αυτή γενικεύεται: χαμηλή ισχύς δεν συνεπάγεται χαμηλή ενέργεια. Η
αρχή ότι η μετακίνηση των δεδομένων κοστίζει ενεργειακά περισσότερο από τον ίδιο
τον υπολογισμό~\cite{chen2016eyeriss} ισχύει και εδώ· εκείνο που λείπει είναι η
επαναχρησιμοποίηση που θα την απόσβενε. Μια
αρχιτεκτονική χαμηλής κατανάλωσης αποδίδει ενεργειακά μόνο αν ολοκληρώνει την
εργασία σε συγκρίσιμο χρόνο, και για το συγκεκριμένο πρόβλημα αυτό δεν συμβαίνει.

\begin{figure}[H]
    \centering
    \includegraphics[width=0.9\textwidth]{int/energy/time_vs_energy}
    \caption[Συνολικός χρόνος και ενέργεια ανά μετρημένη περίπτωση]{Συνολικός χρόνος και συνολική ενέργεια κάθε μετρημένης περίπτωσης.
    Και τα δύο μεγέθη ελαχιστοποιούνται, οπότε οι προτιμότερες περιπτώσεις
    βρίσκονται προς την κάτω αριστερή γωνία.}
    \label{fig:eval-time-energy}
\end{figure}

Το σχήμα~\ref{fig:eval-time-energy} τοποθετεί κάθε μετρημένη περίπτωση σε
επίπεδο δύο κριτηρίων που ελαχιστοποιούνται ταυτόχρονα. Μια περίπτωση λέγεται
ότι κυριαρχείται από μια άλλη όταν η δεύτερη δεν είναι χειρότερη σε κανένα από
τα δύο μεγέθη και είναι καλύτερη σε τουλάχιστον ένα· μια τέτοια περίπτωση δεν
υπάρχει λόγος να επιλεγεί ποτέ. Το σύνολο των περιπτώσεων που δεν κυριαρχούνται
από καμία άλλη αποτελεί το μέτωπο \tl{Pareto}, δηλαδή ακριβώς εκείνες τις
επιλογές στις οποίες η βελτίωση του ενός κριτηρίου απαιτεί υποχώρηση στο άλλο.

Το αποτέλεσμα είναι διαφωτιστικό ακριβώς επειδή είναι αρνητικό. Οι μετρήσεις
σχηματίζουν στενή διαγώνια ζώνη, με συντελεστή συσχέτισης των λογαρίθμων χρόνου
και ενέργειας ίσο με 0,97. Από τις 209 περιπτώσεις μόνο δύο ανήκουν στο
μέτωπο, και αμφότερες αφορούν την ίδια υλοποίηση, το \tl{map-reduce} του
επεξεργαστή, σε δύο γειτονικά μεγέθη εισόδου. Δεν υφίσταται δηλαδή ουσιαστικός
συμβιβασμός μεταξύ ταχύτητας και ενέργειας για το πρόβλημα αυτό: η ταχύτερη
διαμόρφωση είναι ταυτόχρονα και η οικονομικότερη.

Ο λόγος προκύπτει από τα δύο προηγούμενα σχήματα. Η ισχύς μεταβάλλεται μεταξύ
των υλοποιήσεων μόλις κατά παράγοντα 2,3, ενώ ο χρόνος μεταβάλλεται κατά
τέσσερις τάξεις μεγέθους. Όταν ο ένας παράγοντας του γινομένου είναι σχεδόν
σταθερός, το γινόμενο ακολουθεί τον άλλον. Η δυνατότητα ανταλλαγής μεγαλύτερης
διάρκειας με χαμηλότερη ισχύ, η οποία σε άλλα συστήματα αποτελεί πραγματικό
σχεδιαστικό μοχλό, εδώ δεν υπάρχει.

Αξίζει τέλος να σημειωθεί ποιες περιπτώσεις συγκροτούν το μέτωπο: όλες
εκτελούνται στον επεξεργαστή. Καμία διαμόρφωση της μονάδας γραφικών ή της
μονάδας νευρωνικής επεξεργασίας δεν είναι βέλτιστη κατά \tl{Pareto} σε κανένα
σημείο του παραμετρικού χώρου που εξετάστηκε.

\subsection{Γινόμενο ενέργειας-καθυστέρησης}
\label{sec:eval-edp}

\begin{figure}[H]
    \centering
    \includegraphics[width=0.9\textwidth]{int/energy/edp_vs_n_net_k8}
    \caption{Γινόμενο ενέργειας-καθυστέρησης συναρτήσει του μεγέθους εισόδου.}
    \label{fig:eval-edp}
\end{figure}

Το γινόμενο ενέργειας-καθυστέρησης συνοψίζει τα δύο κριτήρια σε ένα και
επιβραβεύει τις υλοποιήσεις που είναι ταυτόχρονα γρήγορες και οικονομικές. Η
κατάταξη που προκύπτει είναι η σαφέστερη ολόκληρου του κεφαλαίου. Το
\tl{map-reduce} στον επεξεργαστή έχει το χαμηλότερο γινόμενο από κάθε άλλη
υλοποίηση, με την αμέσως επόμενη, που είναι το \tl{map-reduce} της μονάδας
γραφικών, να υστερεί κατά παράγοντα είκοσι οκτώ και τη χειρότερη κατά
παράγοντα σαράντα τρία.

Το \tl{map-reduce} στον επεξεργαστή έχει το χαμηλότερο γινόμενο σχεδόν σε
ολόκληρο το εύρος· η μοναδική εξαίρεση είναι το μικρότερο μέγεθος του σχήματος,
όπου η καθαρή ενέργειά του μετριέται μηδενική, κάτω από την ανάλυση του
μετρητή σε αυτή τη διάρκεια, αφήνοντας εκεί τη διτονική του επεξεργαστή
ονομαστικά χαμηλότερη. Στα υπόλοιπα μεγέθη, το μέγεθος του χάσματος προς την
επόμενη υλοποίηση δεν είναι μονότονο: κυμαίνεται μεταξύ περίπου μίας και
δυόμισι τάξεων μεγέθους, καθώς η διτονική και το \tl{map-reduce} του
επεξεργαστή εναλλάσσονται στη δεύτερη θέση ανάλογα με το ακριβές μέγεθος
εισόδου. Η κατάταξη αυτή καθαυτή δεν αμφισβητείται σοβαρά ποτέ· μόνο το
περιθώριο είναι θορυβώδες.

Το αποτέλεσμα αυτό είναι το πιο αντίθετο στη διαίσθηση από όσα παρουσιάζονται
εδώ. Για επιλογή \tl{top-k} με μικρό \tl{k}, η εκφόρτωση σε επιταχυντή δεν
απλώς αποτυγχάνει να αποδώσει ενεργειακά· είναι σαφώς χειρότερη από το να μην
γίνει καθόλου.

\section{Επίδραση του τύπου δεδομένων}
\label{sec:eval-dtypes}

\begin{table}[H]
    \centering
    \begin{tabular}{lrrrrr}
        \hline
        Υλοποίηση & \tl{half} & \tl{int} & \tl{uint} & \tl{float} & \tl{double} \\
        \hline
        \tl{CPU} διτονική & 4,60 & 8,01 & 7,93 & 8,01 & 14,03 \\
        \tl{CPU map-reduce} & 0,39 & 1,23 & 1,22 & 1,22 & 2,35 \\
        \tl{CPU} αναφορά & 13,27 & 4,65 & 4,63 & 5,50 & 5,54 \\
        \tl{GPU} διτονική & 0,88 & 1,34 & 1,37 & 1,34 & 2,89 \\
        \tl{NPU} διτονική & 7,71 & 8,17 & 8,37 & 8,23 & 25,38 \\
        \hline
    \end{tabular}
    \caption[Αλγοριθμικός χρόνος ανά τύπο δεδομένων]{Αλγοριθμικός χρόνος ανά τύπο δεδομένων, σε \tl{ms}, για $2^{24}$ στοιχεία και \tl{k} ίσο με 8.}
    \label{tab:eval-dtypes}
\end{table}

Οι υλοποιήσεις της εργασίας ακολουθούν το πλάτος του τύπου, όπως αναμένεται από
αλγορίθμους περιορισμένους από τη μνήμη. Η μισή ακρίβεια είναι σταθερά η
ταχύτερη, με τη διτονική ταξινόμηση στον επεξεργαστή να ολοκληρώνει σε λίγο
πάνω από τον μισό χρόνο της απλής, ακριβώς επειδή διακινεί τα μισά \tl{bytes} με
τον ίδιο ρυθμό συγκρίσεων.

Η υλοποίηση αναφοράς συμπεριφέρεται αντίθετα: η μισή ακρίβεια είναι η
βραδύτερη περίπτωσή της. Η διαφορά αποκαλύπτει ότι η βιβλιοθήκη δεν περιορίζεται
από τη μνήμη, αλλά από το κόστος της βαθμωτής σύγκρισης, το οποίο για τη μισή
ακρίβεια περιλαμβάνει μετατροπή. Η κατάταξη μεταξύ τύπων δεν είναι επομένως
ιδιότητα του υλικού, αλλά της υλοποίησης.

Η διπλή ακρίβεια στη μονάδα νευρωνικής επεξεργασίας ξεχωρίζει, με τριπλάσιο
χρόνο από κάθε άλλον τύπο. Η επιβάρυνση οφείλεται εξ ολοκλήρου στο πρόσθετο
πέρασμα διόρθωσης της ενότητας~\ref{sec:npu-refinement}, το οποίο απαιτείται
επειδή το κλειδί των 32 \tl{bit} δεν αρκεί για να εκφράσει τιμές των 64.

Ο συνδυασμός των δύο αντίθετων συμπεριφορών παράγει το μοναδικό σημείο ολόκληρης
της εργασίας όπου η μονάδα νευρωνικής επεξεργασίας υπερτερεί της βιβλιοθήκης της.
Για μισή ακρίβεια, η διτονική ταξινόμηση στη μονάδα είναι ταχύτερη από τη
βιβλιοθήκη σε είκοσι εννέα από τις μετρούμενες περιπτώσεις, με τη μεγαλύτερη
διαφορά να φτάνει τον παράγοντα δυόμισι. Το φαινόμενο εμφανίζεται ασθενέστερα και
στους δύο άλλους τύπους κινητής υποδιαστολής, ενώ για τους ακέραιους δεν
εμφανίζεται ποτέ.

Η αιτία δεν είναι ότι η μονάδα χειρίζεται καλά τη μισή ακρίβεια, αλλά ότι η
βιβλιοθήκη τη χειρίζεται άσχημα. Ο πυρήνας της διάταξης λειτουργεί σε ακεραίους
των 32 \tl{bit} ανεξαρτήτως τύπου, οπότε το κόστος του τύπου μετατοπίζεται στην
κωδικοποίηση του \tl{host}, η οποία για τη μισή ακρίβεια είναι διανυσματική
σύμφωνα με την ενότητα~\ref{sec:npu-encoding}. Η βιβλιοθήκη αντιθέτως εκτελεί
βαθμωτές συγκρίσεις με ενδιάμεση μετατροπή. Η εκφόρτωση αποδεικνύεται δηλαδή
συμφέρουσα ακριβώς εκεί όπου η καθιερωμένη λύση είναι ασθενέστερη, και όχι εκεί
όπου η μονάδα είναι ισχυρότερη.

\section{Ευαισθησία στην κατανομή της εισόδου}
\label{sec:eval-distributions}

\begin{figure}[H]
    \centering
    \begin{subfigure}{\textwidth}
        \centering
        \includegraphics[width=\textwidth]{float/dist/dist_sensitivity_k8}
        \caption{\tl{k}=8.}
        \label{fig:eval-dists-k8}
    \end{subfigure}

    \vspace{0.3em}

    \begin{subfigure}{\textwidth}
        \centering
        \includegraphics[width=\textwidth]{float/dist/dist_sensitivity_k131072}
        \caption{\tl{k}=131072.}
        \label{fig:eval-dists-k131072}
    \end{subfigure}
    \caption[Ευαισθησία του χρόνου στην κατανομή της εισόδου]{Ευαισθησία του χρόνου στην κατανομή της εισόδου, για απλή ακρίβεια
    και $2^{24}$ στοιχεία. Κάθε στήλη δίνει τον χρόνο μιας κατανομής ως προς τον
    χρόνο της ομοιόμορφης για την ίδια υλοποίηση· η μονάδα σημαίνει πλήρη
    ανεξαρτησία από τα δεδομένα.}
    \label{fig:eval-dists}
\end{figure}

Η κανονικοποίηση ως προς την ομοιόμορφη κατανομή επιτρέπει να συγκριθούν στο ίδιο
σχήμα και οι εννέα συνδυασμοί, παρότι οι απόλυτοι χρόνοι τους διαφέρουν κατά
τάξεις μεγέθους. Μια υλοποίηση ανεπηρέαστη από τα δεδομένα εμφανίζεται ως πέντε
ίσες στήλες ύψους ένα, ενώ μια ευαίσθητη εμφανίζεται με στήλες που αποκλίνουν
προς τις δύο κατευθύνσεις.

Το σχήμα επιβεβαιώνει πρώτα το αναμενόμενο: η διτονική ταξινόμηση στον
επεξεργαστή και στη μονάδα γραφικών είναι πρακτικά ανεπηρέαστη από την κατανομή
και στις δύο τιμές του \tl{k}, με συνολική διακύμανση κάτω του παράγοντα 1,12. Η
ιδιότητα αυτή την καθιστά αξιόπιστο μάρτυρα: αν ο χρόνος της μεταβαλλόταν, η
μεταβολή θα αποδιδόταν στη μέτρηση και όχι στον αλγόριθμο.

Οι υλοποιήσεις που βασίζονται σε σωρούς εμφανίζουν το αντίθετο άκρο, και μάλιστα
ανεξάρτητα από το ποιος έγραψε τον κώδικα. Η βιβλιοθήκη αναφοράς μεταβάλλεται
κατά παράγοντα εβδομήντα δύο μεταξύ αύξουσας και φθίνουσας εισόδου, αλλά και το
δικό μας \tl{map-reduce} στον επεξεργαστή μεταβάλλεται κατά παράγοντα είκοσι
πέντε, για τον ίδιο μηχανισμό της ενότητας~\ref{sec:cpu-mapreduce}. Η ευαισθησία
στα δεδομένα είναι επομένως ιδιότητα του αλγορίθμου και όχι της υλοποίησης: όποιος
επιλέγει με σωρό την κληρονομεί.

Η βιβλιοθήκη αναφοράς της μονάδας γραφικών εμφανίζει τρίτο μοτίβο:
επιβραδύνεται σχεδόν κατά τον ίδιο παράγοντα και προς τις δύο κατευθύνσεις,
4,95 για την αύξουσα έναντι 4,92 για τη φθίνουσα, αντί να ευνοεί τη μία όπως
οι σωροί. Καθώς η ταξινόμηση δεν μεταβάλλει την κατανομή των τιμών
(ενότητα~\ref{sec:meth-distributions}), αυτό δεν είναι η κατανομή \tl{radix-adversarial}
ψηφίων της ενότητας~\ref{sec:bg-radix}, αλλά ζήτημα θέσης: οι υποψήφιοι που
επιζούν του κρίσιμου κάδου καταλήγουν συγκεντρωμένοι στο ένα ή στο άλλο άκρο
του πίνακα, κάτι που επιβαρύνει συμμετρικά το παράλληλο στάδιο συμπύκνωσης
της βιβλιοθήκης ανεξάρτητα από το για ποιο άκρο πρόκειται.

Δύο παρατηρήσεις αξίζουν ιδιαίτερης προσοχής. Πρώτον, η διτονική ταξινόμηση στη
μονάδα νευρωνικής επεξεργασίας \emph{είναι} ευαίσθητη στην κατανομή, σε αντίθεση
με τις δύο άλλες αρχιτεκτονικές, με διακύμανση παράγοντα πενήντα τρία. Η
ευαισθησία δεν προέρχεται από τη διάταξη, η οποία εκτελεί σταθερό δίκτυο, αλλά
από τη συγχώνευση με σωρό που εκτελεί ο \tl{host} και η οποία, όπως δείχθηκε στην
ενότητα~\ref{sec:npu-pipeline}, κυριαρχεί στον χρόνο.

Δεύτερον, η αντιπαράθεση των δύο τιμών του \tl{k} απομονώνει την αιτία της
ευαισθησίας. Το \tl{map-reduce} στη μονάδα γραφικών είναι ευαίσθητο για \tl{k}
ίσο με 8, όπου η ταξινομημένη είσοδος το επιβραδύνει κατά παράγοντα 4,3, αλλά
γίνεται σχεδόν ανεπηρέαστο για \tl{k} ίσο με 131072, όπου η ίδια είσοδος το
μεταβάλλει κατά μόλις 4\%. Ο ίδιος κώδικας, με τα ίδια δεδομένα, αλλάζει
συμπεριφορά επειδή αλλάζει ο περιοριστικός παράγοντας: στη δεύτερη περίπτωση δεν
τον περιορίζει η απόρριψη, αλλά η έλλειψη παραλληλίας της
ενότητας~\ref{sec:gpu-mapreduce-limit}. Όταν ο περιοριστικός παράγοντας δεν είναι
τα δεδομένα, η κατανομή τους παύει να έχει σημασία.

Οι μετρήσεις που παρουσιάστηκαν καλύπτουν ολόκληρο τον παραμετρικό χώρο της
ενότητας~\ref{sec:meth-parameters} και συμφωνούν μεταξύ τους ως προς το ποιος
παράγοντας περιορίζει κάθε υλοποίηση. Η ερμηνεία τους, τα όρια εντός των οποίων
ισχύουν και οι κατευθύνσεις που ανοίγουν εξετάζονται στο
κεφάλαιο~\ref{ch:conclusions}.
